\documentclass[11pt,a4paper]{article}
\usepackage[margin=25mm]{geometry}
\usepackage{amsmath,amssymb,amsthm,hyperref}
\newtheorem{lemma}{Lemma}\newtheorem{proposition}[lemma]{Proposition}
\newtheorem{theorem}[lemma]{Theorem}
\newcommand{\F}{\mathbb F}\newcommand{\Om}{\Omega}
\newcommand{\Sp}{\operatorname{Sp}}\newcommand{\Aut}{\operatorname{Aut}}
\newcommand{\PCO}{\operatorname{PCO}}\newcommand{\GU}{\operatorname{GU}}
\newcommand{\Tr}{\operatorname{Tr}}
\setlength{\emergencystretch}{2em}
\title{Counting overgroups of an orthogonal torus element}
\author{Human Compression project}\date{}
\begin{document}\maketitle
This working proof uses the previously reviewed orthogonal tori and
fixed-group estimates. The maximal subgroup classifications cited below
are retained inputs, not results newly proved here. The plus eight-dimensional
case has a separate, still unpaid classification interface; the final
paragraphs for that case are conditional on its stated cover.

Write $\nu(H)=|\{h\in H:h^2=1\}|$, so that the identity is included.
For any homomorphism of finite groups with kernel $K$,
\begin{equation}\label{kernel}
                         \nu(H)\le |K|\nu(H/K).
\end{equation}
We use the accepted criterion that an element $x\in G$ has a generating
involution $y$ with no involutory automorphism simultaneously inverting
$x$ and fixing $y$, provided
\begin{equation}\label{criterion}
 i_2(G)> \sum_{\alpha^2=1,\ x^\alpha=x^{-1}}i_2(C_G(\alpha))
             +|N_G(\langle x\rangle)|+\sum_M i_2(M).
\end{equation}
Here the last sum is over the actual maximal subgroups containing $x$.
Replacing these involution counts by $\nu$ only enlarges the right side.
The normalizer is charged even when already contained in that union.

\section{The chosen elements and the classification input}
On a negative space of dimension $2k$ identify the underlying space with
$E=\F_{q^{2k}}$ and use
\[
 Q(v)=\Tr_{\F_{q^k}/\F_q}(c v^{q^k+1}),\qquad c\ne0.
\]
Let $g$ be a primitive norm-one multiplier. The accepted torus proof
establishes its spinor norm, exact projective order, centralizer and
eigenvalue degree. Use the following lifts $X$ of $x$:
\begin{enumerate}
\item $G=\Om_{2m+1}(q)$, $q$ odd, $m\ge3$: $g^2$ on the negative
$2m$-space and the identity on its orthogonal line.
\item $G=P\Om^-_{2m}(q)$, $m\ge4$: $g$ for even $q$, and $g^2$
for odd $q$.
\item $G=P\Om^+_{2m}(q)$, $m\ge4$: $g\oplus h$, with arbitrary
primitive multipliers on negative spaces of dimensions $2m-2$ and two.
\end{enumerate}
The connected projective conformal centralizer has respective orders
\begin{equation}\label{tori}
             |T|=q^m+1,\quad q^m+1,\quad(q^{m-1}+1)(q+1).
\end{equation}

We retain the following complete overgroup types from
\cite[Lemmas~6.1,7.3,8.1]{BSW}. The statements there concern the
corresponding subgroups of $\Om$; lifting a projective maximal subgroup
through the centre gives exactly their applicable setting.
In odd dimension the only type is $\Om^-_{2m}(q).2$, except for
$(m,q)=(3,3)$, which is already in the separately acquired finite list.
In minus dimension the types are the extension field normalizers
\[
       \Om^-_{2m/r}(q^r).r,
       \qquad r\mid m\text{ prime},\quad 2m/r\ge4,
\]
and a unitary normalizer of type $\GU_m(q).2$ when $m$ is odd.
In plus dimension $2m\ge10$ they are the stabilizer of the negative
$(2m-2)+2$ decomposition and, for even $m$, a normalizer of type
$\GU_m(q).2$. Lemma~8.1 uses primitive components of precisely the
orders used here. We do not replace them by their squares.

\section{Counting actual subgroups}
A list of conjugacy types is insufficient for \eqref{criterion}.
We next bound the number of actual subgroups of each type containing
the fixed element.

\begin{lemma}\label{multiplicity}
The odd-dimensional overgroup and the plus decomposition overgroup
are unique. In minus dimension there is at most one extension field
normalizer for each prime $r$, and at most one unitary normalizer.
In plus dimension there are at most two unitary normalizers.
\end{lemma}
\begin{proof}
The distinct primary spaces of $X$ determine the line or the two
negative summands uniquely. This proves the first assertion.

Suppose $X$ lies in an extension field normalizer with scalar field
$K=\F_{q^r}$. Its image in the field automorphism quotient has order
dividing $r$, so $X^r$ is $K$-linear. Its order is at least
$(q^m+1)/(2r)$. Since $r\le m/2$ and $m\ge4$,
\begin{equation}\label{degreebound}
            \frac{q^m+1}{2r}\ge\frac{q^m+1}{m}>q^{m/3}+1.
\end{equation}
For completeness, factor
$q^m+1=(q^{m/3}+1)(q^{2m/3}-q^{m/3}+1)$. The second factor
increases with $q$. For $q=2,m\ge4$, put $a=2^{m/3}>5/2$;
its value $a^2-a+1$ exceeds four at $m=4$, and its derivative
in $m$ is $(\log2)(2a^2-a)/3>1$ throughout the range.
Thus the second factor exceeds $m$, proving \eqref{degreebound}.

The acquired degree lemma for divisors of $q^m+1$ now makes $X^r$
irreducible. Thus $\F_q[X^r]=E$, a field of order $q^{2m}$, and
$K\subseteq C_{\operatorname{End}_{\F_q}(V)}(X^r)=E$. There is only
one subfield of that order. Moreover its transferred quadratic form
is unique. Indeed, if $Q(v)=\Tr_{K/\F_q}(Q_K(v))$, then the values
$Q(av)$ for all $a\in K$ determine $Q_K(v)$ by trace nondegeneracy;
in characteristic two the squaring map on $K$ is bijective, and in odd
characteristic the squares span $K$ additively. Consequently the full
normalizer of this field and form is unique. Taking its intersection
with the fixed ambient group cannot create a second maximal subgroup.

For a minus unitary normalizer, $X^2$ centralizes its quadratic scalar
field. Its order is at least $(q^m+1)/4$ in odd characteristic and the
full order $q^m+1$ in even characteristic. The same acquired degree
lemma, and $(q^m+1)/4>q^{m/3}+1$ for odd $q\ge3,m\ge3$, identify
the scalar field uniquely inside $E$. The hermitian form is recovered
from the ground-field polar form by trace pairing, so its normalizer
is unique too.

For a plus unitary normalizer, $X^2$ is irreducible on the large
summand, by that argument with $m-1$. Its quadratic scalar field is
therefore the unique quadratic subfield of the large endomorphism
field. In particular $X$, which is a multiplier there, centralizes
this field. The field automorphism induced by an element of the
global unitary normalizer is the same on every summand; it must
therefore be trivial everywhere. On the two-space $X$ itself, rather
than $X^2$, determines the quadratic scalar field, since $h$ is
primitive of order $q+1>2$. This covers $q=3$ too. There are two
possible relative identifications of the two quadratic fields,
modulo a simultaneous field automorphism. For each the polar trace
pairing determines the hermitian form. Thus at most two normalizers
occur. These are bounds on actual field structures and forms, not
just their conjugacy classes.
\end{proof}

\begin{lemma}\label{normalizer}
In all three families, $|N_G(\langle x\rangle)|\le4m|T|$.
\end{lemma}
\begin{proof}
Conjugating $x$ to a generator of its cyclic group conjugates $X$ to
a scalar multiple of a coprime power of $X$. It therefore preserves
the algebra generated by $X$, and its primary spaces. On an
irreducible $2k$-space a normalizer of the field algebra has at most
$2k$ possible Frobenius actions, each with fibre its multiplier group.
There is one such block in the first two families, and in plus type
there are at most $2(m-1)$ choices on the large block and two on the
small block. The conformal condition and division by common scalars
give the fibre $T$ of \eqref{tori}. These bounds are at most $4m|T|$.
Only elements of $G$ are counted, so triality automorphisms do not
enlarge this normalizer.
\end{proof}

\begin{lemma}\label{overbound}
The actual overgroups in Lemma~\ref{multiplicity} have total $\nu$
bounded by the applicable terms in the following list:
\begin{equation}\label{overgroupsum}
\begin{array}{ll}
\text{odd-dimensional negative overgroup}&6q^{m^2},\\
\text{plus decomposition}&12(q+1)q^{(m-1)^2},\\
\text{all unitary normalizers}&8(q+1)q^{(m^2+m-2)/2},\\
\text{all minus extension field normalizers}&8m q^{\lfloor m^2/2\rfloor}.
\end{array}
\end{equation}
\end{lemma}
\begin{proof}
The acquired conformal bounds give
$\nu(\PCO_d(q))\le3q^{\lfloor d^2/4\rfloor}$ for odd $q$, and
the corresponding bound with constant four for even $q$.
For even $q$ the latter follows from the retained automorphism
involution estimate \cite[Proposition~1.3]{LLS}, including the
separate small derived-group identifications. In dimension two the
negative conformal group is dihedral, with $\nu=q+3$ for odd $q$
and $\nu=q+2$ for even $q$; both are at most $2(q+1)$.

Restriction to the negative hyperplane in odd total dimension has
projective kernel of order at most two, giving $6q^{m^2}$ by
\eqref{kernel}. The decomposition subgroup maps into the product
of the conformal quotients on its summands; its kernel has order at
most two for odd $q$, and one for even $q$. This gives respectively
$12(q+1)q^{(m-1)^2}$ and the smaller $8(q+1)q^{(m-1)^2}$.

Each unitary normalizer maps to $\Aut(\operatorname{PSU}_m(q))$
with scalar kernel of size at most $q+1$. Its simple derived group
is in the valid range $m\ge4$, including the usual small isomorphisms.
The retained automorphism estimate, with the identity added, is at
most $4q^{(m^2+m-2)/2}$. Lemma~\ref{multiplicity} allows at most
two normalizers, giving the third term.

An extension field orthogonal normalizer acts on the corresponding
derived projective group over $q^r$, with kernel of order at most
two. The same estimate is at most $4q^{m^2/r}$; in dimension four
minus type use $P\Om^-_4(q^r)\cong\operatorname{PSL}_2(q^{2r})$.
The field automorphisms are already included. There are at most
$m$ allowable primes $r$, and $r\ge2$, giving the last bound.
Faithfulness modulo the indicated scalar kernels is the standard
natural-action interface of these classical groups; it is retained
explicitly, rather than proved by a bound on their orders.
\end{proof}

\section{Natural reversing automorphisms}
The torus proof excludes nontrivial ordinary field involutions inverting
$x$. Outside plus dimension eight the remaining diagram automorphisms
are represented on the natural module. Quotients of two such inverters
centralize $x$, so there are at most $|T|$ of them, whether that coset
is inner-diagonal or the other natural orthogonal coset.
The acquired odd-characteristic fixed-group estimates are
\begin{equation}\label{oddfixed}
 \nu(C_G(\alpha))\le
 \begin{cases}
 18q^{m(m+1)/2},&\dim V=2m+1,\\
 24q^{2\lfloor m^2/4\rfloor},&\dim V=2m.
 \end{cases}
\end{equation}
They include conformal lifts of order four and nonsquare multipliers.

For even $q$ in dimension $2m$, the inverter normalizes to an isometry
of rank $m$ with nonalternating induced form. The accepted exact
centralizer sequence has kernel $q^{m(m-1)/2}$ and quotient
$\Sp_{m-1}(q)$ for odd $m$, and kernel $q^{(m^2+m-2)/2}$ and
quotient $\Sp_{m-2}(q)$ for even $m$. Thus
\begin{equation}\label{evenfixed}
 \nu(C_G(\alpha))\le4q^{e_m},\qquad
 e_m=\begin{cases}(3m^2-2m-1)/4,&m\text{ odd},\\
                  (3m^2-4)/4,&m\text{ even}.
       \end{cases}
\end{equation}
The small quotient $\Sp_4(2)$ has $\nu=76$, below the bound used;
it is not treated as simple. For $m=4$, the exact quotient
$\Sp_2(q)$ has $\nu=q^2$, so the sharper bound is $q^{11}$.

\section{The uniform inequalities}
The acquired involution lower bounds are $q^{m^2+m}/4$ in odd
dimension and $q^{m^2-1}/4$ in even dimension. All are strict.

\begin{proposition}\label{uniform}
Criterion~\eqref{criterion} holds in odd dimension $2m+1$ outside
$(m,q)=(3,3)$, and in both even-dimensional families for $m\ge5$,
outside $(m,q)=(5,2),(6,2)$.
\end{proposition}
\begin{proof}
In odd dimension, division by $q^{m^2+m}$ bounds the right side by
\begin{equation}\label{normalizedB}
 18(1+q^{-m})q^{-m(m-1)/2}+6q^{-m}
                 +4m(1+q^{-m})q^{-m^2}.
\end{equation}
For $m=3,q\ge5$ its maximum is at $q=5$ and is less than $194/1000$.
For $m\ge4,q\ge3$ the maximum is at $(4,3)$ and is less than $1/10$.
Each term decreases with $q$ and with $m$ in these ranges: for the
last term the coefficient ratio $(m+1)/m$ is smaller than the lost
factor $q^{2m+1}$. Both maxima are below $1/4$.

For odd $q$ and $m\ge5$, put $D=m^2-1$. We have $|T|\le3q^m/2$.
Using Lemma~\ref{overbound} as an overcount for either sign, the normalized
right side is at most
\begin{equation}\label{normalizedOddD}
 36q^{m+2\lfloor m^2/4\rfloor-D}
 +16q^{m^2-2m+2-D}+11q^{(m^2+m)/2-D}
 +8m q^{\lfloor m^2/2\rfloor-D}+6m q^{m-D}.
\end{equation}
At $(m,q)=(5,3)$ this is less than $25/1000<1/4$.
The successive gaps in the first exponent are $7,11,17,23,\ldots$;
the remaining gaps are $2m-3$, $(m^2-m-2)/2$,
$\lceil m^2/2\rceil-1$, and $m^2-m-1$.
All increase, and the gap increases in the terms with coefficient $m$
more than pay the factor $(m+1)/m$. Thus every term decreases.

For even $q$, use $|T|\le2q^m$. The corresponding upper bound is
\begin{equation}\label{normalizedEvenD}
 8q^{e_m+m-D}+18q^{m^2-2m+2-D}+12q^{(m^2+m)/2-D}
 +8m q^{\lfloor m^2/2\rfloor-D}+8m q^{m-D}.
\end{equation}
The first exponent gap $D-m-e_m$ equals $k^2-1$ at $m=2k+1$
and $k(k-2)$ at $m=2k$. It is three for $m=5,6$, eight for
$m=7,8$, and increases in consecutive odd-even pairs thereafter.
At $(5,4)$ and $(6,4)$ the total in \eqref{normalizedEvenD} is less than $127/1000$.
For $m\ge7,q\ge2$ it is maximized at $(7,2)$, where it is less
than $41/1000$. The other exponent gaps increase as in \eqref{normalizedOddD}, again
paying the coefficients $m$. All these bounds are below $1/4$.
The omitted even-field parameters are in the acquired finite list.
\end{proof}

\section{Dimension eight}
The acquired odd-characteristic subspace counts give
$i_2(P\Om^+_8(q))>q^{16}/4$ and
$i_2(P\Om^-_8(q))>q^{16}/2$.
Here is the even-characteristic substitute.

\begin{lemma}\label{evenlower}
For even $q$, both $\Om^+_8(q)$ and $\Om^-_8(q)$ have more than
$q^{16}/2$ involutions.
\end{lemma}
\begin{proof}
On $\F_q^4\oplus\F_q^4$ use
\[
 Q_b(x,y)=\Bigl(\sum_i x_i\Bigr)^2+x^Ty+\sum_i b_i y_i^2,
 \qquad t(x,y)=(x+y,y).
\]
The polar form is nondegenerate. Its Arf class is that of $\sum b_i$;
choosing its absolute trace to be zero or one gives the two types.
The map $t$ preserves $Q_b$, has rank four, and has Dickson invariant
zero, so belongs to $\Om$. The already proved centralizer formula
gives $|C_{\Om}(t)|=q^9|\Sp_2(q)|$. One actual $\Om$ conjugacy
class therefore has size
\[
 q^2(q^4-1)^2(q^6-1)\quad\text{or}\quad
 q^2(q^8-1)(q^6-1),
\]
respectively. The first ratio to $q^{16}$ is at least
$(15/16)^2(63/64)>1/2$; the second is larger. This counts one
actual class and needs no assertion about fusion of orthogonal classes.
\end{proof}

In minus dimension eight the only extension prime in the cover is
$r=2$, and there is no unitary type. For odd $q\ge3$ the bound
\begin{equation}\label{minusEight}
 24(q^4+1)q^8+8q^8+16(q^4+1)<q^{16}/2
\end{equation}
follows by dividing by $q^{16}$: the resulting decreasing function
is less than $302/1000$ at $q=3$.

For even $q\ge4$, write $T_+=(q^3+1)(q+1)$. The bound
\begin{equation}\label{evenEight}
 q^{11}T_++20(q+1)q^9+32q^8+16T_+<q^{16}/2
\end{equation}
holds since the ratio is decreasing and is less than $324/1000$ at
$q=4$. It pays minus type already; in plus type it pays the natural
decomposition and two unitary normalizers, if they form the complete
cover after excluding the line stabilizers by the triality argument.

More precisely, the following last two bounds are conditional on
that complete plus-dimensional cover. The separately reviewed triality
argument chooses the primitive components so that for odd $q\ge7$
there are no line overgroups in any of the three eight-dimensional
representations, and at most $T_+$ inverters. Then
\begin{equation}\label{plusEight}
 24T_+q^8+20(q+1)q^9+16T_+<q^{16}/4,
\end{equation}
as the decreasing ratio is less than $12/1000$ at $q=7$.
For $q=5$ there are at most two inverter cosets and six actual
$\Om_7(q)$ line stabilizers, each with $\nu\le3q^{12}$.
Consequently
\begin{equation}\label{plusFive}
 48T_+q^8+18q^{12}+20(q+1)q^9+16T_+<q^{16}/4,
\end{equation}
whose ratio at five is less than $124/1000$.
The cases $q=2,3$ in plus type and $q=2$ in minus type have acquired
finite witnesses. No multiplicity-two assumption for the $q=5$ line
overgroups is used. The missing complete-cover classification must
still be supplied before \eqref{evenEight}--\eqref{plusFive} prove the plus dimension-eight
generation theorem.

\paragraph{Attribution.}
The finite classical group orders, trace and form theory, natural
automorphism actions, and the quoted maximal subgroup classifications
are borrowed. The actual-field multiplicity arguments, the consistent
projective centralizers, the new fixed-group counts and the uniform
inequalities are the contributions used here. Full generation outside
the explicitly conditional case is relative to those named inputs;
no newly solved conjecture, additional original, Lean result or numerical
compression factor is claimed by this working note.

\begin{thebibliography}{9}
\bibitem{BSW} D.~Bubboloni, P.~Spiga and Th.~Weigel,
\emph{Normal $2$-coverings of the finite simple groups and their
generalizations}, arXiv:2208.08756v1, Lemmas~6.1,7.3,8.1.
\bibitem{LLS} R.~Lawther, M.~W.~Liebeck and G.~M.~Seitz,
\emph{Fixed point ratios in actions of finite exceptional groups of Lie type},
Pacific J. Math. 205 (2002), 393--464, Proposition~1.3.
\end{thebibliography}
\end{document}
